\section{Podsumowanie}

Przeprowadzona analiza potwierdziła, że zbiór danych wine jest odpowiedni do zadania klasyfikacji z użyciem sieci neuronowej. Dane są kompletne, wszystkie cechy mają charakter numeryczny, a rozkłady części zmiennych różnią się pomiędzy klasami. Jako metodę wstępnego przetwarzania zaproponowano standaryzację cech wejściowych, jako sposób podziału danych przyjęto proporcję 80\%/20\% z zachowaniem udziału klas, natomiast jako model klasyfikacyjny zaproponowano sieć MLP o architekturze $13 \rightarrow 9 \rightarrow 3$, uczoną metodą propagacji wstecznej błędu z wykorzystaniem algorytmu Adam.